The mechanical reengineering and automation of a Cartesian robotic arm were carried out for the packaging line of the Process Design Laboratory at Universidad del Valle de Guatemala. The original guidance, made of metal tubes and printed carriages, showed friction, misalignment and uneven gantry advance that lubrication did not correct. The carriages, the central assembly and the vertical mechanism were redesigned in Autodesk Inventor to incorporate 12~mm steel shafts and linear bearings, and 48 parts were printed in PLA.

Control was divided between a Portenta H7 with Machine Control, which drives the stepper motors and reads the limit switches and the belt encoder, and an ESP32, which handles the wireless controller, the OLED display, the gripper and a HUSKYLENS 2 camera. The link between both controllers was migrated from I2C to UART over RS485 because of persistent packet corruption. The camera locates the parts with a YOLO11n segmentation model and a homography calibrated with four AprilTags; the encoder makes it possible to predict the pick instant while the belt is moving.

In 120 automatic cycles, the part was picked and delivered in 61.7~\% of the cases: 73.3~\% with the belt at 25~\% and 50.0~\% at 35~\%. Vertical parts reached 90.6~\%, horizontal parts 45.2~\% and diagonal parts 54.4~\%. Homing showed a standard deviation of the measured travel of 1.3~mm in X and 0.5~mm in Y. The accuracy and repeatability of the end effector were not measured.
